

\lussec 8. Strategy for the computation of $\mib\ZA$ and ${\mib c}_{\bf fl}$

The PCAC relation (6.32) holds provided 
the relevant renormalization constants and improvement coefficients 
are chosen appropriately. 
As explained below,
the constant $\ZA$ and the 
coefficient $\cfl$ can conversely be determined,
up to terms of order $a^2$ and $a$, respectively,
by requiring the relation to be satisfied exactly at
two points in the space of kinematical parameters.

The computation through numerical simulation 
of $\ZA$ and $\cfl$ along these lines
assumes that the improvement coefficients 
$\csw$ and $\ca$ have already been calculated.
Periodic boundary conditions should be
imposed in the space directions, but apart from this no
further assumptions are made. In particular,
the boundary conditions in time do not need to be
specified, since the
computed values of $\ZA$ and $\cfl$ 
depend on them only
at the level of the lattice-spacing effects.


\lussubsec 8.1 Integrated form of the PCAC relation

The O($a$) improvement unfortunately leads to an inflation
of terms when the renormalized PCAC relation is expanded
in the correlation functions of the unimproved bare fields.
A slight simplification can however be achieved by summing
the equation over an interval of time $x_0$ 
around $y_0$ and over
all points $\vec{x}$ on the spatial lattice. 
One of the bare correlation functions entering the 
relation is then
\equation{
  \CPP{rs}(t,d)=\sum_{x_0=y_0-d}^{y_0+d}\sum_{\vec{x}}
  \langle\Pax^{rs}(x)\Pax_t^{sr}(y)\rangle,
  \enum
}
and there are $6$ further such correlation functions in which
$P^{rs}$ is replaced by
\equation{
  \tilde{P}^{rs},\hat{P}^{rs},Q^{rs},\ring{\partial}_{\mu}A_{\mu}^{rs},
  \ring{\partial}_{\mu}\tilde{A}_{\mu}^{rs}
  \kern0.5em\hbox{or}\kern0.5em
  \drvstar{\mu}\drv{\mu}P^{rs}.
  \enum
}
In the case of the last three fields in this list, 
the sum over $x_0$ reduces to 
a sum of terms at the boundary of the summation range as in 
\equation{
  \CAP{rs}(t,d)=
  \frac{1}{2}\sum_{\vec{x}}\bigl\{
  \langle(\left.A^{rs}_0(x)\right|_{x_0=y_0+d+1}+
          \left.A^{rs}_0(x)\right|_{x_0=y_0+d})
  P^{sr}_t(y)\rangle
  \noenum
  \nexteq{1.0ex}
  {\phantom{\CAP{rs}(t,d)=
  \frac{1}{2}\sum_{\vec{x}}}}\kern-0.65em
  -\langle(\left.A^{rs}_0(x)\right|_{x_0=y_0-d-1}+
           \left.A^{rs}_0(x)\right|_{x_0=y_0-d})
  P^{sr}_t(y)\rangle
  \bigr\},
  \enum
}
for example.

For notational convenience, it is now helpful 
to introduce the abbreviations
\equation{
  \ZAh{rs}=\ZA\left(1+b_A\mq{rs}+\bar{b}_A\tr\Mq\right),
  \enum
  \nexteq{2.5ex}
  \ZPh{rs}=\ZP\left(1+b_P\mq{rs}+\bar{b}_P\tr\Mq\right).
  \enum
}
Defining the bare current-quark mass sums through
\equation{
  m_{rs}={\ZPh{rs}\over\ZAh{rs}}(\mR{r}+\mR{s}),
  \enum
}
the integrated PCAC relation then becomes
\equation{
  \ZAh{rs}\bigl\{\CAP{rs}+
  \tilde{c}_A\CAtP{rs}+c_A\CdPP{rs}
  -m_{rs}\bigl(\CPP{rs}+\tilde{c}_P\CPtP{rs}\bigr)\bigr\}
  \noenum
  \nexteq{2.0ex}
  \qquad+\CPtP{rs}+\hat{c}_A\CAtP{rs}-2\cfl\CPhP{rs}
  -b_{\chi}\frac{1}{2}(\mq{r}-\mq{s})\CQP{rs}=\rmO(a^2).
  \enum
}
This equation still looks fairly complicated, but one should keep in mind
that it is a linear relation among computable correlation functions,
which must hold for all flow times $t>0$, time ranges
$d\geq0$ and quark flavours
$r\neq s$. The unknown coefficients are thus strongly
constrained by the identity.

At values of $d$ larger than the smoothing range 
$\sqrt{8t}$ of the gradient flow, $\CAtP{rs}$ falls off like a Gaussian
and rapidly becomes negligible with respect to the other terms in
eq.~(8.7). Moreover, the sum of the terms in the curly bracket 
as well as all other terms in the equation 
are then practically independent of $d$.
Without further notice, only this 
range of $d$ is considered in the following and the
two terms proportional to $\CAtP{rs}$ are dropped.


\lussubsec 8.2 Up-down flavour channel

In QCD with mass-degenerate up and down quarks,
the PCAC relation in the up-down channel assumes the form
\equation{
  \ZAh{ud}\bigl\{\CAP{ud}+c_A\CdPP{ud}-m_{ud}\CPP{ud}\bigr\}
  -2\cfl\CPhP{ud}
  =-\bigl(1-\ZAh{ud}\tilde{c}_Pm_{ud}\bigr)\CPtP{ud}
  \enum
}
(for simplicity, the O($a^2$) error term is dropped from now on).
The bare mass combination $m_{ud}$ can be determined as usual
through the PCAC relation at vanishing flow time
and may be assumed to be known at this point.
By computing the correlation functions in eq.~(8.8) at 
two values of the flow time $t$, one then obtains
two linear equations for the ratios
\equation{
  {\ZAh{ud}\over1-\ZAh{ud}\tilde{c}_Pm_{ud}}
  \quad\hbox{and}\quad
  {\cfl\over1-\ZAh{ud}\tilde{c}_Pm_{ud}}.
  \enum
}
Clearly, as long as the chosen flow times are
in the scaling range and are held fixed in physical units,
the calculated ratios do not depend on them,
up to terms of order $a^2$
and $a$, respectively. 

At quark masses close to their physical values, and lattice spacings
$a\leq0.1$ fm,
the factor $1-\ZAh{ud}\tilde{c}_Pm_{ud}$ differs
from unity by a fraction of a percent.
An estimate of the size of this correction
term can be obtained by inserting the known value of $m_{ud}$ and 
the tree-level value (6.22) of $\tilde{c}_{P}$.
$\ZAh{ud}$ and $\cfl$ can thus be computed up to this
small correction factor. Note that the bare pion decay constant
determined through the vacuum-to-pion matrix element of the improved
axial current is to be renormalized with the
mass-dependent factor $\ZAh{ud}$ rather than $\ZA$. The difference
is however small and can be estimated using the 
tree-level value $b_A=1$.


\lussubsec 8.3 Up-strange flavour channel

The additional term that appears in the PCAC relation
\equation{
  \ZAh{us}\bigl\{\CAP{us}+c_A\CdPP{us}-m_{us}\CPP{us}\bigr\}
  -2\cfl\CPhP{us}-b_{\chi}\frac{1}{2}(\mq{u}-\mq{s})\CQP{us}
  \noenum
  \nexteq{2.0ex}
  \qquad
  =-\bigl(1-\ZAh{us}\tilde{c}_Pm_{us}\bigr)\CPtP{us}
  \enum
}
in this channel is proportional to the square of the mass 
difference $\mq{u}-\mq{s}$ and is therefore 
likely to be negligible in practice.
To be able to determine the ratios (8.9) (with $ud$ replaced by $us$),
eq.~(8.10) must otherwise be evaluated at
three values of the flow time $t$.

For $a\leq0.1$ fm and physical quark masses,
the correction factor $1-\ZAh{us}\tilde{c}_Pm_{us}$ 
may deviate from unity by up to one percent or so.
The comparison with the results obtained in the $ud$ channel,
\equation{
  {\ZAh{us}(1-\ZAh{ud}\tilde{c}_Pm_{ud})
  \over\ZAh{ud}(1-\ZAh{us}\tilde{c}_Pm_{us})}
  =1+b_A\frac{1}{2}(\mq{s}-\mq{d})+\ZA\tilde{c}_P(m_{us}-m_{ud}),
  \enum
}
unfortunately only allows a certain linear combination of the coefficients 
$b_A$ and $\tilde{c}_P$ to be estimated.
